\section{Logistic Regression}\label{sec:lr}

\subsection{Giới thiệu về Logistic Regression}
Logistic Regression \cite{hosmer2013} ước lượng xác suất rời bỏ bằng cách đưa một tổ hợp tuyến tính của đặc trưng qua hàm sigmoid; tương đương, log-odds là hàm tuyến tính của đặc trưng, nên hệ số $w_j$ là mức thay đổi log-odds khi $x_j$ tăng một đơn vị:
\begin{equation}
    P(y = 1 \mid \mathbf{x}) = \frac{1}{1 + e^{-(\mathbf{w}^\top \mathbf{x} + b)}}, \qquad
    \log \frac{P(y=1 \mid \mathbf{x})}{1 - P(y=1 \mid \mathbf{x})} = \mathbf{w}^\top \mathbf{x} + b
\end{equation}
Tham số được học bằng cách cực tiểu hóa log-loss có trọng số lớp $s_i$ và chính quy hóa L2 ($C$ nhỏ: chính quy hóa mạnh):
\begin{equation}
    \mathcal{L}(\mathbf{w}, b) = -\sum_{i=1}^{n} s_i \left[y_i \log \hat{p}_i + (1-y_i)\log(1-\hat{p}_i)\right] + \frac{1}{2C}\|\mathbf{w}\|_2^2
\end{equation}
Với \texttt{class\_weight="balanced"}, mẫu rời bỏ có trọng số $s_i \approx 1{,}88$ và mẫu không rời bỏ $s_i \approx 0{,}68$. Hàm mất mát lồi nên có nghiệm tối ưu toàn cục (tìm bằng \texttt{lbfgs}).

\begin{table}[H]
\centering
\begin{adjustbox}{max width=\textwidth}
\begin{tabular}{|p{7cm}|p{7cm}|}
\hline
\textbf{Ưu điểm} & \textbf{Nhược điểm} \\
\hline
Nhanh, ổn định, nghiệm tối ưu toàn cục, ít overfit & Chỉ biểu diễn được biên quyết định tuyến tính \\
\hline
Hệ số và tỉ số odds dễ diễn giải & Không tự học được tương tác, quan hệ phi tuyến \\
\hline
Là baseline mạnh trên dữ liệu bảng & Cần chuẩn hóa; hệ số khó đọc khi đặc trưng tương quan \\
\hline
\end{tabular}
\end{adjustbox}
\caption{Ưu và nhược điểm của Logistic Regression}
\end{table}

\subsection{Mã nguồn và huấn luyện}
Logistic Regression là mô hình duy nhất được bật \texttt{StandardScaler} cho ba đặc trưng số (\texttt{requires\_scaling()}); các cột one-hot giữ nguyên giá trị 0/1. Mô hình được khởi tạo trong \texttt{src/models.py}:

\begin{lstlisting}[style=pythonstyle, caption={Khởi tạo Logistic Regression trong \texttt{src/models.py}}]
class_weight = "balanced" if balance_ratio else None

if model_name == "logistic_regression":
    return LogisticRegression(
        solver="lbfgs",
        max_iter=2000,
        class_weight=class_weight,
        random_state=random_state,
    )
\end{lstlisting}

Siêu tham số duy nhất được tối ưu là $C$, lấy mẫu theo phân phối log-uniform trên $[10^{-3}, 10^{2}]$ với 30 tổ hợp (\texttt{src/tune.py}). Giá trị tốt nhất theo PR-AUC là $C = 0{,}92$, gần với mặc định $C = 1$.

\subsection{Kết quả và đánh giá}

\subsubsection{Kết quả cross-validation qua các giai đoạn}
\begin{table}[H]
\centering
\small
\begin{adjustbox}{max width=\textwidth}
\begin{tabular}{|l|c|c|c|c|c|c|}
\hline
\textbf{Cấu hình} & \textbf{Precision} & \textbf{Recall} & \textbf{F1} & \textbf{ROC-AUC} & \textbf{PR-AUC} & \textbf{F1 train} \\
\hline
Baseline & 0,652 & 0,543 & 0,592 & 0,846 & 0,662 & 0,602 \\
\hline
+ Class weight & 0,517 & 0,801 & 0,628 & 0,846 & 0,660 & 0,634 \\
\hline
+ Tuning ($C = 0{,}92$) & 0,517 & 0,801 & 0,628 & 0,846 & 0,660 & 0,634 \\
\hline
+ Ngưỡng 0,60 (out-of-fold) & 0,571 & 0,719 & \textbf{0,637} & 0,846 & 0,660 & -- \\
\hline
\end{tabular}
\end{adjustbox}
\caption{Logistic Regression: kết quả validation (trung bình 5 fold) qua từng bước}
\end{table}

Class weight là bước mang lại cải thiện chính (F1 +0,036, Recall từ 0,543 lên 0,801, tăng ở cả 5 fold); tuning gần như không thay đổi gì vì mô hình vốn không overfit (F1 train và validation chênh nhau dưới 0,01).

\subsubsection{Kết quả trên tập test}
\begin{table}[H]
\centering
\begin{minipage}[t]{0.56\textwidth}
\centering
\small
\begin{adjustbox}{max width=\textwidth}
\begin{tabular}{|l|c|c|c|c|}
\hline
\textbf{Lớp} & \textbf{Precision} & \textbf{Recall} & \textbf{F1-score} & \textbf{Support} \\
\hline
Không rời bỏ (0) & 0,88 & 0,78 & 0,83 & 1.035 \\
\hline
Rời bỏ (1) & 0,54 & 0,71 & 0,61 & 374 \\
\hline
Accuracy & & & 0,76 & 1.409 \\
\hline
Macro avg & 0,71 & 0,74 & 0,72 & 1.409 \\
\hline
Weighted avg & 0,79 & 0,76 & 0,77 & 1.409 \\
\hline
\end{tabular}
\end{adjustbox}
\caption{Logistic Regression: kết quả phân loại trên tập test (ngưỡng 0,60; ROC-AUC 0,842; PR-AUC 0,633)}
\end{minipage}\hfill
\begin{minipage}[t]{0.42\textwidth}
\centering
\small
\begin{adjustbox}{max width=\textwidth}
\begin{tabular}{|c|c|c|c|}
\cline{3-4}
\multicolumn{2}{c|}{} & \multicolumn{2}{c|}{\textbf{Dự đoán}} \\
\cline{3-4}
\multicolumn{2}{c|}{} & Không rời bỏ & Rời bỏ \\
\hline
\multirow{2}{*}{\textbf{Thực tế}} & Không rời bỏ & \cellcolor{blue!30}809 & 226 \\
\cline{2-4}
 & Rời bỏ & 110 & \cellcolor{blue!30}264 \\
\hline
\end{tabular}
\end{adjustbox}
\caption{Logistic Regression: Confusion Matrix trên tập test}
\end{minipage}
\end{table}

Mô hình phát hiện được 264/374 khách hàng rời bỏ, bỏ sót 110; đổi lại có 226/1.035 khách hàng ở lại bị cảnh báo nhầm.

\subsubsection{Hệ số mô hình}
Vì đặc trưng số đã được chuẩn hóa, hệ số của chúng tương ứng với mức thay đổi log-odds khi đặc trưng tăng một độ lệch chuẩn; với cột one-hot, hệ số là mức thay đổi khi khách hàng thuộc nhóm đó.

\begin{table}[H]
\centering
\small
\begin{adjustbox}{max width=\textwidth}
\begin{tabular}{|l|c|c|l|c|c|}
\hline
\multicolumn{3}{|c|}{\textbf{Tăng khả năng rời bỏ}} & \multicolumn{3}{c|}{\textbf{Giảm khả năng rời bỏ}} \\
\cline{1-3} \cline{4-6}
\textbf{Đặc trưng} & $\beta$ & $e^{\beta}$ & \textbf{Đặc trưng} & $\beta$ & $e^{\beta}$ \\
\hline
InternetService = Fiber optic & 0,686 & 1,98 & \texttt{tenure} & $-$1,151 & 0,32 \\
\hline
Contract = Month-to-month & 0,656 & 1,93 & Contract = Two year & $-$0,774 & 0,46 \\
\hline
\texttt{TotalCharges} & 0,484 & 1,62 & \texttt{MonthlyCharges} & $-$0,632 & 0,53 \\
\hline
StreamingMovies = Yes & 0,265 & 1,30 & InternetService = DSL & $-$0,609 & 0,54 \\
\hline
StreamingTV = Yes & 0,251 & 1,29 & InternetService = No & $-$0,268 & 0,77 \\
\hline
PaymentMethod = Electronic check & 0,240 & 1,27 & & & \\
\hline
OnlineSecurity = No & 0,199 & 1,22 & & & \\
\hline
TechSupport = No & 0,165 & 1,18 & & & \\
\hline
\end{tabular}
\end{adjustbox}
\caption{Các hệ số lớn nhất của Logistic Regression đã tuning}
\end{table}

Permutation importance trên tập test (Hình~\ref{fig:perm}, Mục~\ref{sec:xgb}) xếp \texttt{tenure} quan trọng nhất (PR-AUC giảm 0,232 khi xáo trộn), tiếp theo là \texttt{InternetService} (0,084), \texttt{MonthlyCharges} (0,069) và \texttt{Contract} (0,056).

\subsubsection{Nhận xét và đánh giá}
\textbf{Điểm mạnh}
\begin{itemize}
    \item Là mô hình tốt nhất ở cấu hình baseline (F1 0,592, ROC-AUC 0,846) và gần như không overfit; điểm test khớp với cross-validation (ROC-AUC 0,846 so với 0,842).
    \item Chỉ cần class weight là đạt gần mức của các mô hình phức tạp hơn sau khi tuning (F1 validation 0,628--0,637).
    \item Hệ số dễ diễn giải và nhất quán với EDA: \texttt{tenure} dài và hợp đồng 2 năm làm giảm mạnh nguy cơ rời bỏ (tỉ số odds 0,32 và 0,46); Fiber optic và hợp đồng theo tháng làm odds tăng gần gấp đôi.
    \item Huấn luyện rất nhanh (khoảng 0,04 giây mỗi fold, ngang Decision Tree, so với 0,37 giây của Random Forest 100 cây).
\end{itemize}

\textbf{Hạn chế}
\begin{itemize}
    \item F1 test thấp nhất trong bốn mô hình sau khi tối ưu (0,611), và PR-AUC giảm nhiều nhất từ CV sang test (0,660 xuống 0,633).
    \item Không biểu diễn được quan hệ phi tuyến, ví dụ nguy cơ rời bỏ giảm rất nhanh trong năm đầu rồi đi ngang.
    \item \textbf{Đa cộng tuyến:} \texttt{TotalCharges} xấp xỉ \texttt{tenure} $\times$ \texttt{MonthlyCharges}, còn \texttt{MonthlyCharges} gần như được xác định bởi các dịch vụ đăng ký. Vì vậy dấu của từng hệ số riêng lẻ (\texttt{MonthlyCharges} âm, \texttt{TotalCharges} dương) không nên diễn giải tách rời.
\end{itemize}

\textbf{Hướng cải thiện}
\begin{enumerate}
    \item Thêm đặc trưng phi tuyến cho mô hình tuyến tính (spline hoặc chia nhóm \texttt{tenure}, tương tác hợp đồng $\times$ thời gian sử dụng).
    \item Thử chính quy hóa L1 hoặc Elastic Net để giảm ảnh hưởng của các đặc trưng tương quan và chọn lọc đặc trưng.
    \item Loại bỏ hoặc biến đổi \texttt{TotalCharges} để hệ số dễ diễn giải hơn.
\end{enumerate}
